\documentclass[12pt]{extarticle}
\usepackage{amsmath}
\usepackage{amsfonts}
\usepackage{textcomp}
\usepackage{cite}
\usepackage{multicol}
\usepackage{setspace}
\usepackage{changepage}

\newtheorem{thm}{Theorem}[section]
\newtheorem{cor}{Correlation}[thm]
\newtheorem{defn}{Definition}[section]
\newtheorem{prop}{Preposition}[section]
\newtheorem{lmm}{Lemma}[section]
\newtheorem{conj}{Conjecture}[section]
\newtheorem{exm}{Example}[section]
\newtheorem{rem}{Remark}[section]
\newtheorem{note}{Note}[section]
\newtheorem{alg}{Algorithm}[section]

\setcounter{thm}{1}
\setcounter{cor}{1}
\setcounter{defn}{1}
\setcounter{prop}{1}
\setcounter{lmm}{1}
\setcounter{conj}{1}
\setcounter{exm}{1}
\setcounter{rem}{1}
\setcounter{note}{1}
\setcounter{alg}{1}

\begin{document}
\section{Introduction}

Recently, I realized that it is of immense difficulty to keep too many open tabs in my mind while doing research. Besides, I believe providing people with better context regarding my research is a rewarding activity, specially when they decide to attend my presentations. Therefore, the idea of a project to keep track of the context of the papers and/or books I study, was born.

These notes are available publically under MIT license.

P.S: \emph{You may find some music recommendations within the text. Click on the link and enjoy!}    
\section{Preliminaries}
This section contains the prerequisites to have a better understanding of the papers.
\subsection{Category Theory}
\begin{defn}[Category\cite{Awodey2010-AWOCTA}]
    A \emph{category} consists of the following data:
    \begin{enumerate}
        \item \emph{Objects}: \(A, B, C,\dots\)
        \item \emph{Arrows}: \(f, g, h,\dots\)
        \item For each arrow \(f : A \rightarrow B\), there are given objects \(dom(f)\) \emph{[read: "domain of f"]} and \(cod(f)\) \emph{[read: "codomain of f"]}, respectively, and namely \(A\) and \(B\) in the example above.
        \item Given arrows \(f: A \rightarrow B\) and \(g : B \rightarrow C\) \([dom(f) = cod(f)]\), there is a given arrow \(g \circ f\ : A \rightarrow C \), called the \emph{composition of \(f\) and \(g\)}.
        \item For each object \(A\), there is a given arrow \(1_{A} : A \rightarrow A\), namely the identity arrow of \(A\).
    \end{enumerate}
    which satisfy the following rules:
    \begin{itemize}
        \item \emph{Assoiativity}: \( \forall f : A \rightarrow B, g: B \rightarrow C, h: C \rightarrow D, \ h \circ (g \circ f) = (h \circ g) \circ f\)
        \item \emph{Unit}: \( \forall f : A \rightarrow B, \ f \circ 1_A = f = 1_B \circ f\)
    \end{itemize}
\end{defn}

Mappings between the categories are called \emph{functors}. They map one category to another while preserving its sturcture. 

\begin{defn}[Functors]
    A \emph{functor} \(F : C \rightarrow D\) between categories \(\mathbf{C}\) and \(\mathbf{D}\) is a mapping of objects to objects and arrows to arrows in such a way that:
    \begin{enumerate}
        \item \(F (f: A \rightarrow B) = F(f) : F(A) \rightarrow F(B)\)
        \item \(F(1_A) = 1_{F(A)}\)
        \item \(F(g \circ f) = F(g) \circ F(f)\)
    \end{enumerate}
\end{defn}

\subsection{Nominal Sets}

\section{Papers}

\bibliographystyle{abbrv}
\bibliography{ref}
\end{document}